% Luc Destombe --- licence CC BY-NC-SA 4.0
% https://creativecommons.org/licenses/by-nc-sa/4.0/deed.fr
% Fichier autonome : compiler avec pdflatex, deux fois (signets, nombre de pages).
\documentclass[a4paper,10pt]{article}
\makeatletter
\newif\iflycee@niveaux
\lycee@niveauxtrue
\newif\iflycee@palatino
\lycee@palatinotrue

\RequirePackage[utf8]{inputenc}
\RequirePackage[T1]{fontenc}
\RequirePackage[french]{babel}
\RequirePackage{etoolbox}

\newcommand{\ShorthandsOff}{\@ifundefined{shorthandoff}{}{\shorthandoff{;:!?}}}
\newcommand{\ShorthandsOn}{\@ifundefined{shorthandon}{}{\shorthandon{;:!?}}}
\ShorthandsOff
\AtBeginDocument{\ShorthandsOn}
\AtBeginEnvironment{tikzpicture}{\ShorthandsOff}

\RequirePackage{geometry}
\RequirePackage{fancyhdr}
\RequirePackage{lastpage}
\RequirePackage{multicol}
\RequirePackage{array}
\RequirePackage{enumitem}
\RequirePackage{xcolor}
\RequirePackage{needspace}

\RequirePackage{amsmath,amssymb}
\RequirePackage{mathpazo}
\DeclareMathSizes{10.95}{10.95}{8}{5.5}
\begingroup\catcode`\_=\active \catcode`\^=\active
\gdef_#1{\sb{\scriptscriptstyle #1}}
\gdef^#1{\sp{\scriptscriptstyle\lycee@moinsepais #1}}
\endgroup
\newcommand\lycee@moins{\mathbin{%
  \setbox\z@\hbox{$\m@th\scriptscriptstyle\lycee@moinsorig$}%
  \dimen@=.55\wd\z@ \advance\dimen@-.5pt
  \hbox to.75\wd\z@{\hss$\m@th\scriptscriptstyle\vcenter{\hbox{%
    \tikz\draw[line width=.5pt,line cap=round](0,0)--(\the\dimen@,0);}}$\hss}}}
\begingroup\lccode`\~=`\- \lowercase{\endgroup\let~}\lycee@moins
\newcommand\lycee@moinsepais{\mathcode`\-="8000 }
\AtBeginDocument{\mathchardef\lycee@moinsorig=\mathcode`\-\relax}
\AtBeginDocument{\catcode`\_=12 \mathcode`\_="8000
  \catcode`\^=12 \mathcode`\^="8000 }
\newcommand\lycee@exposants{%
  \fontdimen13\textfont2=.48\fontdimen6\textfont2
  \fontdimen14\textfont2=.48\fontdimen6\textfont2
  \fontdimen15\textfont2=.48\fontdimen6\textfont2
  \fontdimen13\scriptfont2=.48\fontdimen6\scriptfont2
  \fontdimen14\scriptfont2=.48\fontdimen6\scriptfont2
  \fontdimen15\scriptfont2=.48\fontdimen6\scriptfont2}
\every@math@size\expandafter{\the\every@math@size\lycee@exposants}
\newlength{\retraitcalculs}
\setlength{\retraitcalculs}{2em}

\RequirePackage{tikz}
\usetikzlibrary{arrows.meta,calc,babel,shapes.misc}
\RequirePackage[most]{tcolorbox}
\tcbuselibrary{skins,breakable}

\definecolor{coulDef}{HTML}{1F4E79}
\definecolor{coulProp}{HTML}{7030A0}
\definecolor{coulRem}{HTML}{C55A11}
\definecolor{coulEx}{HTML}{2E7D32}
\definecolor{coulExo}{HTML}{B03A2E}
\definecolor{coulPart}{HTML}{1F4E79}
\definecolor{coulMeth}{HTML}{1B6E4F}
\definecolor{coulCourbe}{HTML}{0B5ED7}
\definecolor{coulCourbeB}{HTML}{C00000}
\definecolor{coulCourbeC}{HTML}{2E7D32}
\definecolor{coulGrille}{HTML}{8C8C8C}
\definecolor{coulRep}{HTML}{1B6E4F}
\definecolor{coulBar}{HTML}{2E7D32}

\newcommand{\R}{\mathbb{R}}
\newcommand{\Rs}{\mathbb{R}^{*}}
\renewcommand{\le}{\leqslant}
\renewcommand{\ge}{\geqslant}
\newcommand{\vect}[1]{\overrightarrow{#1}}
\newcommand{\norme}[1]{\left\lVert #1 \right\rVert}
\newcommand{\intervalle}[2]{\left[#1\,;#2\right]}
\RequirePackage{cancel}
\renewcommand{\CancelColor}{\color{coulExo}}
\newcommand{\fois}[1]{\times{\color{coulExo}#1}}
\newcommand{\barre}[1]{\cancel{#1}}

\tikzset{grille/.style={coulGrille,line width=0.4pt}}
\newcommand{\quadrillage}[4]{\draw[grille,step=1] (#1,#2) grid (#3,#4);}
\tikzset{
  croix/.style={cross out,draw,line width=0.6pt,inner sep=0pt,minimum size=4pt},
  croixdroite/.style={croix,rotate=45,line width=0.8pt},
}
\newcommand{\pt}[3]{\path (#1) node[croix]{} node[#2,font=\small,inner sep=2.5pt]{$#3$};}

\newcommand{\lycee@repere}[4]{%
  \draw[grille,step=1] (#1,#2) grid (#3,#4);
  \draw[-{Stealth[length=2mm]},thick] (#1,0)--({#3+0.4},0) node[below left=-1pt]{$x$};
  \draw[-{Stealth[length=2mm]},thick] (0,#2)--(0,{#4+0.4}) node[below left=-1pt]{$y$};
  \node[below left,font=\scriptsize,inner sep=1.5pt] at (0,0) {$0$};}
\newcommand{\lycee@gradx}[1]{\foreach \k/\l in {#1}{%
  \draw (\k,0.07)--(\k,-0.07) node[below,font=\scriptsize,inner sep=1.5pt]{$\l$};}}
\newcommand{\lycee@grady}[1]{\foreach \k/\l in {#1}{%
  \draw (0.07,\k)--(-0.07,\k) node[left,font=\scriptsize,inner sep=1.5pt]{$\l$};}}
\newcommand{\lycee@intff}[2]{\left[#1\,;#2\right]}
\newcommand{\lycee@intfo}[2]{\left[#1\,;#2\right[}
\newcommand{\lycee@intof}[2]{\left]#1\,;#2\right]}
\newcommand{\lycee@intoo}[2]{\left]#1\,;#2\right[}
\newcommand{\lycee@ens}[1]{\left\{#1\right\}}
\AtBeginDocument{%
  \providecommand{\repere}{\lycee@repere}%
  \providecommand{\gradx}{\lycee@gradx}%
  \providecommand{\grady}{\lycee@grady}%
  \providecommand{\Cf}{\mathcal{C}\sb{\scriptscriptstyle f}}%
  \providecommand{\Cg}{\mathcal{C}\sb{\scriptscriptstyle g}}%
  \providecommand{\intff}{\lycee@intff}%
  \providecommand{\intfo}{\lycee@intfo}%
  \providecommand{\intof}{\lycee@intof}%
  \providecommand{\intoo}{\lycee@intoo}%
  \providecommand{\ens}{\lycee@ens}}

\newlist{questions}{enumerate}{3}
\setlist[questions,1]{label=\arabic*\textdegree),leftmargin=*,itemsep=2pt,topsep=3pt}
\setlist[questions,2]{label=\alph*),leftmargin=*,itemsep=1pt,topsep=2pt}
\setlist[questions,3]{label=\roman*),leftmargin=*,itemsep=1pt,topsep=1pt}
\setlist[itemize]{label=\textbullet,leftmargin=*,itemsep=2pt,topsep=3pt}

\newcommand{\besoinplace}[1]{\par\penalty\@M\begingroup
  \dimen@=#1\relax\dimen@ii=\pagegoal\advance\dimen@ii-\pagetotal
  \ifdim\dimen@>\dimen@ii\ifdim\dimen@ii>\z@\vfil\fi\break\fi\endgroup}

\newcommand{\lycee@pied}{\thepage/\pageref{LastPage}}
\newcommand{\entete}[2]{%
  \pagestyle{fancy}\fancyhf{}%
  \fancyhead[L]{\footnotesize\color{coulPart}\textbf{#1}}%
  \fancyhead[R]{\footnotesize\color{coulPart}\textbf{#2}}%
  \fancyfoot[C]{\footnotesize\lycee@pied}%
  \renewcommand{\headrulewidth}{0.4pt}%
  \renewcommand{\headrule}{\hbox to\headwidth{%
    \color{coulPart!40}\leaders\hrule height \headrulewidth\hfill}}}

\RequirePackage{ccicons}
\newcommand{\lycee@licence}{{\scriptsize\color{gray}%
  \copyright\ Luc Destombe --- \ccbyncsa\ CC BY-NC-SA 4.0}}
\newif\iflycee@licence \lycee@licencetrue
\newcommand{\sanslicence}{\lycee@licencefalse}
\AtBeginDocument{\iflycee@licence\AddToHookNext{shipout/foreground}{%
  \put(0,0){\rlap{\kern\dimexpr1in+\hoffset+\oddsidemargin+\textwidth\relax
    \llap{\raisebox{-\dimexpr1in+\voffset+\topmargin+\headheight+\headsep
      +\textheight+\footskip\relax}{\lycee@licence}}}}}\fi}

\newcommand{\formulesserrees}{%
  \setlength{\abovedisplayskip}{3pt}\setlength{\belowdisplayskip}{3pt}%
  \setlength{\abovedisplayshortskip}{2pt}\setlength{\belowdisplayshortskip}{3pt}}

\setlength{\parindent}{0pt}

\RequirePackage[implicit=false,hidelinks,pdfencoding=auto,psdextra,bookmarksopen,
  bookmarksopenlevel=1,pdfcreator={},pdfproducer={}]{hyperref}
\RequirePackage{bookmark}
\hypersetup{pdfauthor={Luc Destombe},pdfkeywords={CC BY-NC-SA 4.0}}
\newcounter{lycee@signet}
\newcommand{\signet}[3][]{\stepcounter{lycee@signet}%
  \edef\lycee@dest{\if\relax\detokenize{#1}\relax
    signet.\arabic{lycee@signet}\else#1\fi}%
  \hypertarget{\lycee@dest}{}\bookmark[level=#2,dest=\lycee@dest]{#3}}
\pdfstringdefDisableCommands{%
  \def\R{R}\def\vect#1{#1}\def\Cf{Cf}\def\Cg{Cg}\def\fois#1{×#1}%
  \def\barre#1{#1}\def\intervalle#1#2{[#1;#2]}\def\intff#1#2{[#1;#2]}%
  \def\textsuperscript#1{#1}\def\dfrac#1#2{#1/#2}\def\frac#1#2{#1/#2}%
  \def\vec#1{#1⃗}}%
\RequirePackage[official]{eurosym}

\geometry{top=0.8cm,bottom=1.3cm,left=1cm,right=1cm,footskip=0.6cm}

\pagestyle{fancy}
\fancyhf{}
\renewcommand{\headrulewidth}{0pt}
\fancyfoot[C]{\footnotesize Page \thepage{} / \pageref{LastPage}}

\newcommand{\titrefiche}[2]{%
  \begin{center}
    \Large\textbf{#1}\\[2pt]
    \large\textbf{#2}\\[2pt]
  \end{center}
  \vspace{0.10cm}}

\setlength{\columnseprule}{0.4pt}
\setlength{\columnsep}{15pt}
\renewcommand{\columnseprulecolor}{\color{black!45}}
\setlength{\multicolsep}{2pt}

\newcounter{partiefiche}
\renewcommand{\thepartiefiche}{\Roman{partiefiche}}
\newif\ifapresbandeau
\newcommand{\lycee@bandeau}[1]{%
  \par\penalty-600
  \vspace{0.08cm}%
  \noindent\colorbox{coulPart}{%
    \parbox{\dimexpr\linewidth-2\fboxsep\relax}{%
      \color{white}\bfseries\raggedright\strut#1\strut}}%
  \nopagebreak\par\nopagebreak
  \vspace{0.06cm}\nopagebreak
  \apresbandeautrue}
\newcommand{\partiefiche}[1]{%
  \refstepcounter{partiefiche}%
  \lycee@bandeau{\signet{1}{\thepartiefiche{} --- #1}\thepartiefiche{} --- #1}}
\newcommand{\partiefichesansnum}[1]{\lycee@bandeau{\signet{1}{#1}#1}}

\newcounter{exo}
\newlength{\espaceexo}\setlength{\espaceexo}{7pt}
\newlength{\espacetitre}\setlength{\espacetitre}{2pt}
\newcommand{\exo}[1][\relax]{%
  \par
  \ifapresbandeau\apresbandeaufalse\else\penalty-150\addvspace{\espaceexo}\fi
  \refstepcounter{exo}%
  \noindent\signet[exercice-\arabic{exo}]{2}{Exercice \arabic{exo}}%
  \textbf{\color{coulExo}\underline{Exercice \theexo}}%
  \ifx\relax#1\relax\else\textbf{ : #1}\fi
  \par\nopagebreak\vspace{\espacetitre}\nopagebreak
  \ignorespaces}

\setlist[enumerate]{nosep,leftmargin=*,itemsep=2pt,topsep=2pt}
\setlist[itemize]{nosep,leftmargin=*,topsep=2pt}
\newcommand{\choix}[3]{\par\hspace*{1em}%
  \makebox[0.3\linewidth][l]{a) #1}\makebox[0.3\linewidth][l]{b) #2}c) #3}
\newcommand{\deux}[2]{\par\hspace*{0.6em}%
  \makebox[0.48\linewidth][l]{#1}#2}
\newcommand{\trois}[3]{\par\hspace*{0.6em}%
  \makebox[0.32\linewidth][l]{#1}\makebox[0.32\linewidth][l]{#2}#3}

  \renewcommand{\exo}[1][\relax]{%
    \ifapresbandeau\apresbandeaufalse\else\par\penalty-150\fi
    \refstepcounter{exo}%
    \vspace{0.05cm}%
    \noindent\signet[exercice-\arabic{exo}]{2}{Exercice \arabic{exo}}%
    \textbf{\color{coulExo}\underline{Exercice \theexo}}%
    \ifx\relax#1\relax\else\textbf{ : #1}\fi%
    \par\nopagebreak
    \ignorespaces}
  \newcommand{\rappel}[1]{%
    \par\vspace{0.06cm}\nopagebreak
    \noindent\fcolorbox{black!35}{black!5}{%
      \parbox{\dimexpr\linewidth-2\fboxsep-2\fboxrule\relax}{%
        \small\textbf{Rappel.} #1}}%
    \par\vspace{0.06cm}\nopagebreak}
  \setlist[enumerate]{nosep,leftmargin=*}
  \setlist[itemize]{nosep,leftmargin=*}
  \setlength{\multicolsep}{3pt plus 1pt minus 1pt}
  \setlength{\premulticols}{10pt}
  \setlength{\postmulticols}{10pt}
  \newlist{expr}{enumerate}{1}
  \setlist[expr]{label={\Alph*\ $=$},nosep,leftmargin=*,itemsep=0pt}

\renewenvironment{center}{\par\addvspace{3pt}\centering}{\par\addvspace{3pt}}
\formulesserrees
\AtBeginDocument{\formulesserrees}
\clubpenalty=10000
\widowpenalty=10000
\displaywidowpenalty=10000
\brokenpenalty=10000
\setlength{\parskip}{0pt}

\RequirePackage{ragged2e}
\setlength{\RaggedRightRightskip}{0pt plus 1fil}
\AtBeginDocument{\RaggedRight\binoppenalty=10000 \relpenalty=10000 }
\g@addto@macro\@parboxrestore{\RaggedRight}
\makeatother
\geometry{top=0.8cm,bottom=1.3cm,left=1cm,right=1cm,footskip=0.7cm,headheight=12pt}

\raggedcolumns
\newcommand{\qcm}[4]{\par\makebox[0.24\linewidth][l]{a) #1}%
  \makebox[0.24\linewidth][l]{b) #2}\makebox[0.24\linewidth][l]{c) #3}d) #4}
\newcommand{\point}[2]{\path[coulCourbe] (#1,#2) node[croix=2.2pt]{};}
\newcommand{\axes}[4]{%
  \draw[->,>=Stealth,black!70] (0,0) -- (#1+0.6,0) node[below,font=\footnotesize] {$n$};
  \draw[->,>=Stealth,black!70] (0,0) -- (0,#2+0.6) node[left,font=\footnotesize] {$u_{n}$};
  \foreach \i in {#3} {\node[below,font=\tiny] at (\i,0) {$\i$};}
  \foreach \i in {#4} {\node[left,font=\tiny] at (0,\i) {$\i$};}
  \node[below left,font=\tiny] at (0,0) {$0$};}

\begin{document}

\titrefiche{1\textsuperscript{re} mathématiques spécifiques --- Suites numériques}%
  {Fiche d'exercices du chapitre}

\begin{multicols}{2}

\partiefiche{Notion de suite}

\exo[Calculer les premiers termes]
Pour chaque suite, calculer les cinq premiers termes.
\trois{$u_{n}=4n-5$}{$v_{n}=n^{2}+1$}{$w_{n}=3^{n}$}

\exo[Calculer un terme]
Calculer les termes demandés.
\begin{enumerate}[label=\alph*)]
  \item $u_{n}=6n+5$ : calculer $u_{9}$.
  \item $v_{n}=3n^{2}-2n+1$ : calculer $v_{4}$.
  \item $w_{n}=\dfrac{2n}{n+1}$ : calculer $w_{0}$, $w_{1}$, $w_{2}$ et $w_{3}$.
  \item $t_{n}=20-n^{2}$ : calculer $t_{3}$.
\end{enumerate}

\exo[Des suites dans la vie courante]
Pour chaque situation, exprimer $u_{n}$ en fonction de $n$, puis calculer le terme
demandé.
\begin{enumerate}[label=\alph*)]
  \item Une piscine vend une carte annuelle à $30$~€, puis chaque entrée $3{,}50$~€.
        $u_{n}$ est le prix payé pour $n$ entrées dans l'année. Calculer $u_{12}$.
  \item Une trottinette en libre-service coûte $1$~€ au départ, puis $0{,}25$~€ par
        minute. $u_{n}$ est le prix d'un trajet de $n$ minutes. Calculer $u_{40}$.
  \item La réserve d'eau d'un village, de $9\,000$~m$^{3}$, baisse de $600$~m$^{3}$
        par mois. $u_{n}$ est la réserve au bout de $n$ mois. Calculer $u_{6}$.
\end{enumerate}

\exo[Ne pas confondre]
Soit la suite définie par $u_{n}=2n+5$.
\begin{enumerate}[label=\alph*)]
  \item Calculer $u_{4}$, puis $u_{4}+1$, puis $u_{5}$.
  \item Exprimer $u_{n+1}$ en fonction de $n$.
\end{enumerate}

\exo[Retrouver le rang]
Soit la suite définie par $u_{n}=80-4n$.
\begin{enumerate}[label=\alph*)]
  \item Calculer $u_{0}$, $u_{1}$ et $u_{10}$.
  \item Pour quel rang $n$ a-t-on $u_{n}=48$ ? et $u_{n}=0$ ?
\end{enumerate}

\partiefiche{Suites définies par récurrence}

\exo[Calculer les termes pas à pas]
Pour chaque suite, calculer les termes de rang $1$ à $4$.
\begin{enumerate}[label=\alph*)]
  \item $u_{0}=2$ et $u_{n+1}=3u_{n}-1$
  \item $v_{0}=5$ et $v_{n+1}=v_{n}+4$
  \item $w_{0}=48$ et $w_{n+1}=\dfrac{w_{n}}{2}$
  \item $t_{0}=20$ et $t_{n+1}=t_{n}-6$
\end{enumerate}

\exo[Avec des nombres négatifs]
Calculer les termes demandés.
\begin{enumerate}[label=\alph*)]
  \item $u_{0}=-1$ et $u_{n+1}=3-u_{n}^{2}$ : calculer $u_{1}$, $u_{2}$ et $u_{3}$.
  \item $v_{0}=1$ et $v_{n+1}=-3v_{n}+2$ : calculer $v_{1}$, $v_{2}$ et $v_{3}$.
\end{enumerate}

\exo[Traduire un énoncé]
Pour chaque situation, donner le premier terme et la relation de récurrence, puis
calculer les deux termes suivants.
\begin{enumerate}[label=\alph*)]
  \item Une tirelire contient $150$~€ ; chaque semaine, on y ajoute $20$~€. On note
        $u_{n}$ la somme après $n$ semaines.
  \item Une algue couvre $3$~m$^{2}$ d'un étang, et sa surface triple chaque
        semaine. On note $a_{n}$ la surface, en m$^{2}$, après $n$ semaines.
  \item Une médiathèque compte $400$ abonnés ; chaque année, $30$ partent et $50$
        arrivent. On note $c_{n}$ le nombre d'abonnés après $n$ années.
\end{enumerate}

\exo[Retrouver la relation]
Pour chaque liste, donner le premier terme et une relation de récurrence qui la
produit.
\trois{a) $4$ ; $9$ ; $14$ ; $19$}{b) $2$ ; $8$ ; $32$ ; $128$}{c) $1$ ; $4$ ; $13$ ; $40$}

\exo[Revenir en arrière]
On sait que $u_{n+1}=u_{n}-6$ et que $u_{3}=11$. Calculer $u_{2}$, $u_{1}$, puis
$u_{0}$.

\exo[Au tableur]
Les rangs $n$ sont en colonne A, les termes en colonne B, le premier terme en B2.
Pour chaque suite, quelle formule saisir en B3, à recopier vers le bas ?
\begin{enumerate}[label=\alph*)]
  \item $u_{0}=30$ et $u_{n+1}=u_{n}+15$.
  \item $v_{0}=200$ et $v_{n+1}=0{,}9v_{n}+30$.
  \item $w_{n}=5n-1$ (la formule utilise la colonne A).
\end{enumerate}

\partiefiche{Représentation graphique d'une suite}

\exo[Représenter une suite]
Pour chaque suite, dresser le tableau de valeurs pour $n$ allant de $0$ à $6$, puis
représenter la suite dans un repère.
\deux{a) $u_{n}=9-2n$}{b) $v_{0}=4$ et $v_{n+1}=2-v_{n}$}

\exo[Lire un graphique]
\begin{minipage}[t]{0.5\linewidth}
\vspace{0pt}
On a représenté les premiers termes d'une suite $(u_{n})$.
\begin{enumerate}[label=\alph*)]
  \item Lire $u_{0}$ et $u_{4}$.
  \item Pour quel rang a-t-on $u_{n}=4$ ?
  \item Les termes augmentent-ils ou diminuent-ils ?
\end{enumerate}
\end{minipage}\hfill
\begin{minipage}[t]{0.47\linewidth}
\vspace{0pt}
\centering
\begin{tikzpicture}[x=0.45cm,y=0.45cm]
  \quadrillage{0}{0}{6}{7}
  \axes{6}{7}{1,2,3,4,5}{1,2,3,4,5,6,7}
  \point{0}{1}\point{1}{2}\point{2}{2.5}\point{3}{4}\point{4}{5.5}\point{5}{7}
\end{tikzpicture}
\end{minipage}

\exo[Calculer puis représenter]
Soit la suite définie par $u_{0}=-6$ et $u_{n+1}=\dfrac{u_{n}+10}{2}$.
\begin{enumerate}[label=\alph*)]
  \item Calculer $u_{1}$, $u_{2}$, $u_{3}$ et $u_{4}$.
  \item Représenter ces cinq termes dans un repère.
\end{enumerate}

\exo[Une suite qui descend puis remonte]
Soit la suite définie par $w_{n}=n^{2}-6n+8$.
\begin{enumerate}[label=\alph*)]
  \item Dresser le tableau de valeurs pour $n$ allant de $0$ à $6$.
  \item Représenter ces termes dans un repère.
  \item Décrire l'évolution des termes.
\end{enumerate}

\partiefiche{Suites arithmétiques : définition}

\exo[Reconnaître une suite arithmétique]
Ces nombres sont-ils les premiers termes d'une suite arithmétique ? Si oui, donner
la raison.
\trois{a) $6$ ; $10$ ; $14$ ; $18$}{b) $12$ ; $9$ ; $6$ ; $3$}{c) $2$ ; $6$ ; $18$ ; $54$}
\deux{d) $1{,}5$ ; $2$ ; $2{,}5$ ; $3$}{e) $0$ ; $1$ ; $4$ ; $9$}

\exo[Modéliser]
Pour chaque situation, donner le premier terme, la relation de récurrence et la
raison.
\begin{enumerate}[label=\alph*)]
  \item Un cinéma vend $50\,000$ places la première année ($u_{0}$), puis $4\,000$ de
        plus chaque année.
  \item Le stock d'un entrepôt, de $6\,000$ cartons, baisse de $450$ cartons par
        semaine.
  \item Une cagnotte contient $200$~€ ; on y ajoute $15$~€ chaque mois.
\end{enumerate}

\exo[Arithmétique ou non ?]
Dire si chaque suite est arithmétique, en justifiant.
\begin{enumerate}[label=\alph*)]
  \item $u_{n}=7-3n$ : calculer $u_{n+1}-u_{n}$. La suite est-elle arithmétique ?
  \item $v_{n}=n^{2}+3$ : calculer $v_{0}$, $v_{1}$ et $v_{2}$. La suite est-elle
        arithmétique ?
  \item $w_{n}=3(n+2)$ : montrer que la suite est arithmétique.
\end{enumerate}

\exo[Compléter]
Les suites $(u_{n})$ et $(v_{n})$ sont arithmétiques.
\begin{enumerate}[label=\alph*)]
  \item $u_{0}=9$ et $r=-4$ : calculer $u_{1}$, $u_{2}$, $u_{3}$ et $u_{4}$.
  \item $v_{2}=13$ et $v_{3}=18$ : donner la raison, puis $v_{4}$, $v_{1}$ et $v_{0}$.
\end{enumerate}

\partiefiche{Suites arithmétiques : expression en fonction de $n$}

\exo[Appliquer la formule]
Pour chaque suite arithmétique, exprimer $u_{n}$ en fonction de $n$, puis calculer
le terme demandé.
\begin{enumerate}[label=\alph*)]
  \item $u_{0}=10$ et $r=-3$ : $u_{8}$.
  \item $u_{0}=-4$ et $r=5$ : $u_{12}$.
  \item $u_{1}=600$ et $r=40$ : $u_{8}$.
  \item $u_{1}=-7$ et $r=4$ : $u_{5}$.
\end{enumerate}

\exo[Une suite qui diminue]
On donne $u_{1}=185$ ; $u_{2}=172$ ; $u_{3}=159$ et $u_{4}=146$.
\begin{enumerate}[label=\alph*)]
  \item Montrer que ce sont les termes d'une suite arithmétique, et donner sa
        raison.
  \item Calculer $u_{5}$, puis $u_{14}$.
\end{enumerate}

\exo[Production d'un atelier]
\label{ex:production}%
En 2010, un atelier fabrique $4\,000$ vélos ; sa production augmente de $250$ vélos
par an. On note $u_{n}$ la production en $2010+n$.
\begin{enumerate}[label=\alph*)]
  \item Donner la nature de la suite $(u_{n})$, puis exprimer $u_{n}$ en fonction de
        $n$.
  \item Calculer la production en 2013, puis en 2030.
\end{enumerate}

\exo[Lire la formule]
Donner le premier terme $u_{0}$ et la raison de chaque suite arithmétique.
\trois{a) $u_{n}=5+4n$}{b) $v_{n}=18-1{,}5n$}{c) $w_{n}=6n$}

\exo[Un terme lointain]
Calculer le terme demandé de chaque suite arithmétique.
\begin{enumerate}[label=\alph*)]
  \item $u_{0}=15$ et $r=4$ : $u_{100}$.
  \item $u_{1}=7$ et $r=3$ : $u_{50}$.
\end{enumerate}

\partiefiche{Suites arithmétiques : problème de seuil}

\exo[Acheter un vélo électrique]
Une lycéenne veut acheter un vélo électrique à $1\,500$~€. Le 1\textsuperscript{er}
janvier, elle a $900$~€, et elle met $60$~€ de côté chaque mois. On note $u_{1}=900$
la somme au 1\textsuperscript{er} janvier, $u_{2}$ celle au 1\textsuperscript{er}
février, et ainsi de suite.
\begin{enumerate}[label=\alph*)]
  \item Calculer $u_{2}$, $u_{3}$ et $u_{4}$.
  \item Exprimer $u_{n}$ en fonction de $n$.
  \item Calculer $u_{7}$.
  \item Pour quelle valeur de $n$ aura-t-elle $1\,500$~€ ?
  \item À quelle date pourra-t-elle acheter son vélo ?
\end{enumerate}

\exo[Doubler la production]
On reprend l'atelier de l'exercice~\ref{ex:production} : $u_{n}=4\,000+250n$. En
quelle année la production aura-t-elle doublé par rapport à 2010 ?

\exo[Le recul d'une falaise]
En 2020, une maison se trouve à $200$~m du bord d'une falaise, qui recule de $2$~m
par an. On note $d_{n}$ la distance, en mètres, en $2020+n$.
\begin{enumerate}[label=\alph*)]
  \item Exprimer $d_{n}$ en fonction de $n$.
  \item À partir de quelle année la maison sera-t-elle à moins de $175$~m du bord ?
\end{enumerate}

\exo[Résoudre des inéquations]
Pour chaque inéquation, donner le plus petit entier $n$ qui la vérifie.
\trois{a) $15+5n>100$}{b) $240-30n<0$}{c) $40+6n\ge 100$}

\partiefiche{Suites arithmétiques : variations et représentation graphique}

\exo[Sens de variation]
Donner la raison, puis le sens de variation de chaque suite arithmétique.
\deux{a) $u_{n}=3+8n$}{b) $v_{n}=20-0{,}4n$}
\par\hspace*{0.6em}c) $t_{n}=-5n$
\par\hspace*{0.6em}d) $w_{0}=-6$ et $w_{n+1}=w_{n}+0{,}5$

\exo[Lire un graphique]
\begin{minipage}[t]{0.5\linewidth}
\vspace{0pt}
Les points ci-contre représentent une suite arithmétique $(u_{n})$.
\begin{enumerate}[label=\alph*)]
  \item Lire $u_{0}$ et la raison $r$.
  \item Exprimer $u_{n}$ en fonction de $n$.
  \item Calculer $u_{10}$.
  \item Donner le sens de variation.
\end{enumerate}
\end{minipage}\hfill
\begin{minipage}[t]{0.47\linewidth}
\vspace{0pt}
\centering
\begin{tikzpicture}[x=0.45cm,y=0.45cm]
  \quadrillage{0}{0}{5}{8}
  \axes{5}{8}{1,2,3,4}{1,2,3,4,5,6,7,8}
  \point{0}{8}\point{1}{6.5}\point{2}{5}\point{3}{3.5}\point{4}{2}
\end{tikzpicture}
\end{minipage}

\exo[Représenter]
Représenter les cinq premiers termes des suites $(u_{n})$ et $(v_{n})$ dans un même
repère, puis donner leur sens de variation.
\deux{a) $u_{n}=1+1{,}5n$}{b) $v_{0}=8$ et $v_{n+1}=v_{n}-2$}

\partiefiche{Suites géométriques : définition}

\exo[Coefficient multiplicateur]
Donner le coefficient multiplicateur $q$.
\deux{a) Augmenter de $4\,\%$}{b) Diminuer de $4\,\%$}
\deux{c) Augmenter de $30\,\%$}{d) Diminuer de $30\,\%$}
\deux{e) Augmenter de $60\,\%$}{f) Diminuer de $25\,\%$}
\deux{g) Augmenter de $2\,\%$}{h) Diminuer de $45\,\%$}

\exo[Retrouver l'évolution]
Traduire chaque coefficient par « augmenter de\ldots{} \% » ou « diminuer de\ldots{}
\% ».
\trois{a) $q=0{,}96$}{b) $q=1{,}06$}{c) $q=0{,}4$}
\trois{d) $q=1{,}3$}{e) $q=3$}{f) $q=0{,}98$}
\trois{g) $q=0{,}85$}{h) $q=1{,}075$}{}

\exo[Reconnaître une suite géométrique]
Ces nombres sont-ils les premiers termes d'une suite géométrique ? Si oui, donner
la raison.
\deux{a) $3$ ; $12$ ; $48$ ; $192$}{b) $64$ ; $32$ ; $16$ ; $8$}
\deux{c) $4$ ; $8$ ; $12$ ; $16$}{d) $500$ ; $550$ ; $605$ ; $665{,}5$}

\exo[Modéliser]
Pour chaque situation, donner le premier terme, la relation de récurrence et la
raison, puis calculer le terme de rang $1$.
\begin{enumerate}[label=\alph*)]
  \item Un capital de $5\,000$~€ est placé à $2\,\%$ par an.
  \item Une forêt de $40\,000$ arbres perd $5\,\%$ de ses arbres par an.
  \item Les émissions de CO$_{2}$ d'une usine, de $800$~t, doivent baisser de
        $4\,\%$ par an.
  \item Une population de $300$ insectes augmente de $50\,\%$ chaque semaine.
\end{enumerate}

\exo[Appliquer une évolution]
Calculer, sans calculatrice.
\deux{a) $300$ augmenté de $20\,\%$}{b) $500$ diminué de $10\,\%$}
\deux{c) $2\,000$ augmenté de $3\,\%$}{d) $60$ diminué de $25\,\%$}

\partiefiche{Suites géométriques : expression en fonction de $n$}

\exo[Appliquer la formule]
Pour chaque suite géométrique, exprimer $u_{n}$ en fonction de $n$, puis calculer le
terme demandé.
\begin{enumerate}[label=\alph*)]
  \item $u_{0}=2$ et $q=3$ : $u_{4}$.
  \item $u_{0}=5\,000$ et $q=0{,}1$ : $u_{3}$.
  \item $u_{1}=6$ et $q=2$ : $u_{5}$.
  \item $u_{0}=128$ et $q=0{,}5$ : $u_{4}$.
\end{enumerate}

\exo[Placer ses économies]
Une étudiante place $1\,500$~€ au taux annuel de $2\,\%$. On note $c_{n}$ le capital
au bout de $n$ années.
\begin{enumerate}[label=\alph*)]
  \item Calculer $c_{1}$ et $c_{2}$.
  \item Donner la nature de $(c_{n})$, puis exprimer $c_{n}$ en fonction de $n$.
  \item On donne $1{,}02^{10}\approx 1{,}22$. Estimer le capital au bout de $10$ ans.
\end{enumerate}

\exo[Covoiturage]
En 2024, $80$ salariés d'une entreprise viennent au travail en covoiturage. On
suppose que ce nombre augmente de $10\,\%$ par an, et on note $u_{n}$ ce nombre en
$2024+n$.
\begin{enumerate}[label=\alph*)]
  \item Donner la nature de la suite $(u_{n})$ et sa raison.
  \item Les années sont en ligne 1 (2024 en B1) et les termes en ligne 2 ($80$ en
        B2). Quelle formule saisir en C2, à recopier vers la droite ?
  \item Exprimer $u_{n}$ en fonction de $n$.
  \item On donne $1{,}1^{3}=1{,}331$. Estimer $u_{3}$ à l'unité près.
  \item L'entreprise vise un quart de ses $400$ salariés en covoiturage en 2027.
        L'objectif sera-t-il atteint ?
\end{enumerate}

\exo[Dioxyde d'azote]
\label{ex:azote}%
Dans une ville, la concentration en dioxyde d'azote est de $60$~µg/m$^{3}$ en 2020 ;
on veut la faire baisser de $5\,\%$ par an. On note $u_{n}$ la concentration en
$2020+n$.
\begin{enumerate}[label=\alph*)]
  \item Calculer $u_{1}$, puis exprimer $u_{n}$ en fonction de $n$.
  \item On donne $0{,}95^{6}\approx 0{,}74$. Estimer la concentration en 2026.
\end{enumerate}

\exo[Lire la formule]
Donner le premier terme $u_{0}$ et la raison de chaque suite géométrique.
\deux{a) $u_{n}=4\times 5^{n}$}{b) $v_{n}=2\,000\times 0{,}85^{n}$}
\par\hspace*{0.6em}c) $w_{n}=3^{n}$

\partiefiche{Suites géométriques : variations et représentation graphique}

\exo[Sens de variation]
Donner la raison, puis le sens de variation de chaque suite géométrique.
\deux{a) $u_{n}=8\times 1{,}3^{n}$}{b) $v_{n}=500\times 0{,}6^{n}$}
\par\hspace*{0.6em}c) $w_{0}=7$ et $w_{n+1}=1{,}02\,w_{n}$
\par\hspace*{0.6em}d) $t_{0}=40$ et $t_{n+1}=0{,}25\,t_{n}$
\par\hspace*{0.6em}e) une quantité qui baisse de $5\,\%$ par an.

\exo[Linéaire ou exponentielle ?]
\begin{minipage}[t]{0.5\linewidth}
\vspace{0pt}
On a représenté deux suites : $(a_{n})$ par des croix droites ($+$), $(b_{n})$ par
des croix obliques ($\times$).
\begin{enumerate}[label=\alph*)]
  \item Laquelle est arithmétique ? géométrique ? Donner leurs raisons.
  \item Calculer $a_{5}$ et $b_{5}$.
  \item Laquelle a une décroissance linéaire ? exponentielle ?
\end{enumerate}
\end{minipage}\hfill
\begin{minipage}[t]{0.47\linewidth}
\vspace{0pt}
\centering
\begin{tikzpicture}[x=0.45cm,y=0.45cm]
  \quadrillage{0}{0}{5}{8}
  \axes{5}{8}{1,2,3,4}{1,2,3,4,5,6,7,8}
  \foreach \n/\v in {0/8,1/6.5,2/5,3/3.5,4/2}
    {\path[coulCourbeB] (\n,\v) node[croixdroite]{};}
  \point{0}{8}\point{1}{4}\point{2}{2}\point{3}{1}\point{4}{0.5}
\end{tikzpicture}
\end{minipage}

\exo[Vrai ou faux ?]
Justifier chaque réponse.
\begin{enumerate}[label=\alph*)]
  \item Une suite géométrique de raison $1{,}05$ (et de premier terme positif) est
        décroissante.
  \item Une quantité qui baisse de $4\,\%$ par an est modélisée par une suite
        décroissante.
  \item Les points d'une suite arithmétique de raison $3$ sont alignés.
  \item $u_{n}=5\times 0{,}8^{n}$ est croissante, car $5>1$.
\end{enumerate}

\partiefiche{Suites géométriques : problème de seuil}

\exo[Dioxyde d'azote (suite)]
On reprend $u_{n}=60\times 0{,}95^{n}$ (exercice~\ref{ex:azote}). La ville vise
$45$~µg/m$^{3}$ au plus.
\begin{enumerate}[label=\alph*)]
  \item Montrer que $u_{n}\le 45$ revient à $0{,}95^{n}\le 0{,}75$.
  \item On donne $0{,}95^{5}\approx 0{,}77$ ; $0{,}95^{6}\approx 0{,}74$ et
        $0{,}95^{7}\approx 0{,}70$. En quelle année l'objectif sera-t-il atteint ?
\end{enumerate}

\exo[Lettre d'information]
En janvier, une lettre d'information compte $2\,500$ abonnés ; ce nombre augmente
ensuite de $6\,\%$ par mois. On note $u_{n}$ le nombre d'abonnés $n$ mois après
janvier.
\begin{enumerate}[label=\alph*)]
  \item Exprimer $u_{n}$ en fonction de $n$.
  \item Montrer que $u_{n}\ge 5\,000$ revient à $1{,}06^{n}\ge 2$.
  \item On donne $1{,}06^{11}\approx 1{,}90$ ; $1{,}06^{12}\approx 2{,}01$ et
        $1{,}06^{13}\approx 2{,}13$. Au bout de combien de mois le nombre d'abonnés
        aura-t-il doublé ?
\end{enumerate}

\exo[La valeur d'un scooter]
Un scooter acheté $3\,000$~€ perd $20\,\%$ de sa valeur chaque année.
\begin{enumerate}[label=\alph*)]
  \item Exprimer sa valeur $v_{n}$ après $n$ années en fonction de $n$.
  \item On donne $0{,}8^{3}\approx 0{,}51$ ; $0{,}8^{4}\approx 0{,}41$ et
        $0{,}8^{5}\approx 0{,}33$. Au bout de combien d'années vaudra-t-il moins de
        $1\,200$~€ ?
\end{enumerate}

\partiefiche{Comparer deux évolutions}

\exo[Arithmétique, géométrique ou aucune des deux ?]
Donner la nature de chaque suite et, s'il y a lieu, sa raison.
\deux{a) $u_{n}=5n$}{b) $u_{n}=n^{2}-1$}
\deux{c) $u_{n}=4\times 3^{n}$}{d) $u_{n+1}=u_{n}-2$}
\deux{e) $u_{n+1}=u_{n}+0{,}2\,u_{n}$}{f) $u_{n+1}=2u_{n}$}

\exo[Deux placements]
En 2025, Yanis dispose de $15\,000$~€ pour un projet qui en coûte $18\,000$.
Placement A : $400$~€ d'intérêts chaque année ; on note $a_{n}$ le capital en
$2025+n$. Placement B : $+3\,\%$ par an ; on note $b_{n}$ le capital en $2025+n$.
\begin{enumerate}[label=\alph*)]
  \item Calculer $a_{1}$ et $b_{1}$.
  \item Donner la nature et la raison de chaque suite.
  \item Exprimer $a_{n}$ et $b_{n}$ en fonction de $n$.
  \item Avec le placement A, en quelle année le projet sera-t-il financé ?
  \item On donne $b_{6}\approx 17\,911$ et $b_{7}\approx 18\,448$. Même question avec
        le placement B. Quel placement choisir ?
\end{enumerate}

\exo[Réduire ses déchets]
Deux entreprises veulent diviser par deux leurs déchets plastiques de 2020.
L'entreprise A en produit $360$~t en 2020, et $15$~t de moins chaque année ;
l'entreprise B en produit $500$~t, et $7\,\%$ de moins chaque année. On note $a_{n}$
et $b_{n}$ leurs productions en $2020+n$.
\begin{enumerate}[label=\alph*)]
  \item Donner la nature et la raison de chaque suite.
  \item Exprimer $a_{n}$ et $b_{n}$ en fonction de $n$.
  \item En quelle année l'entreprise A atteindra-t-elle son objectif ?
  \item On donne $0{,}93^{9}\approx 0{,}52$ et $0{,}93^{10}\approx 0{,}48$. En quelle
        année l'entreprise B atteindra-t-elle son objectif ?
  \item Laquelle y parvient la première ?
\end{enumerate}

\partiefichesansnum{Exercices type épreuve anticipée}

\exo[Automatismes (QCM)]
Pour chaque question, une seule réponse est exacte.
\begin{enumerate}[label=\arabic*.]
  \item Diminuer de $30\,\%$ revient à multiplier par :
        \qcm{$0{,}3$}{$1{,}3$}{$0{,}7$}{$-30$}
  \item La liste $4$ ; $8$ ; $16$ ; $32$ est celle d'une suite :
        \par a) arithmétique de raison $4$ \quad b) géométrique de raison $2$
        \quad c) ni l'une ni l'autre
  \item $u_{0}=5$ et $u_{n+1}=2u_{n}$ ; alors $u_{3}$ vaut :
        \qcm{$10$}{$40$}{$30$}{$20$}
  \item $u_{n}=60-3n$ ; alors $u_{10}$ vaut :
        \qcm{$57$}{$30$}{$600$}{$90$}
  \item Une suite géométrique de raison $1{,}02$ et de premier terme $50$ est :
        \par a) croissante \quad b) décroissante \quad c) constante
  \item $u_{n}=400\times 0{,}5^{n}$ ; alors $u_{1}$ vaut :
        \qcm{$200$}{$400{,}5$}{$800$}{$0{,}5$}
\end{enumerate}

\exo[Location de vélos]
Une jeune entreprise de location possède $40$ vélos au mois $0$ et en achète $30$
chaque mois. Elle compte $25$ clients au mois $0$, puis $20\,\%$ de clients de plus
chaque mois. On note $u_{n}$ le nombre de vélos et $v_{n}$ le nombre de clients au
mois $n$.
\begin{enumerate}[label=\arabic*.]
  \item Exprimer $u_{n}$ en fonction de $n$. Montrer qu'au mois $10$, l'entreprise
        possède plus de $300$ vélos.
  \item Au bout de combien de mois aura-t-elle $640$ vélos ?
  \item Justifier que $v_{n+1}=1{,}2\,v_{n}$. En déduire la nature de $(v_{n})$.
  \item Exprimer $v_{n}$ en fonction de $n$.
  \item On donne $1{,}2^{5}\approx 2{,}49$. Montrer que $v_{5}\approx 62$ et
        interpréter.
  \item On donne $1{,}2^{9}\approx 5{,}16$ et $1{,}2^{10}\approx 6{,}19$. À partir de
        quel mois l'entreprise aura-t-elle plus de $150$ clients ?
\end{enumerate}

\exo[Deux clubs sportifs]
En 2025, un club de tennis compte $700$ adhérents et en perd $15$ chaque année ; on
note $T_{n}$ le nombre d'adhérents en $2025+n$. Un club d'escalade compte $150$
adhérents en 2025 ; le graphique donne ses effectifs $E_{n}$.

\begin{center}
\begin{tikzpicture}[x=0.6cm,y=0.005cm]
  \draw[grille,xstep=1,ystep=100] (0,0) grid (9.5,850);
  \draw[->,>=Stealth,black!70] (0,0) -- (10.1,0) node[below,font=\footnotesize] {$n$};
  \draw[->,>=Stealth,black!70] (0,0) -- (0,950) node[left,font=\footnotesize] {$E_{n}$};
  \foreach \i in {1,...,9} {\node[below,font=\tiny] at (\i,0) {$\i$};}
  \foreach \v in {200,400,...,800} {\node[left,font=\tiny] at (0,\v) {$\v$};}
  \node[below left,font=\tiny] at (0,0) {$0$};
  \foreach \n/\v in {0/150,1/180,2/216,3/259.2,4/311,5/373.2,6/447.9,7/537.5,8/645,9/774}
    {\path[coulCourbe] (\n,\v) node[croix=2.2pt]{};}
\end{tikzpicture}
\end{center}
\begin{enumerate}[label=\arabic*.]
  \item Calculer $T_{1}$, puis exprimer $T_{n}$ en fonction de $n$.
  \item Justifier que le club de tennis aura perdu plus de $20\,\%$ de ses adhérents
        entre 2025 et 2035.
  \item Estimer par lecture graphique le nombre d'adhérents du club d'escalade en
        2028.
  \item On admet que $E_{n+1}=1{,}2\,E_{n}$. Donner la nature de $(E_{n})$, son premier
        terme et sa raison.
  \item On donne $1{,}2^{7}\approx 3{,}58$ et $1{,}2^{8}\approx 4{,}30$. À partir de
        quelle année le club d'escalade aura-t-il plus d'adhérents que le club de
        tennis ?
\end{enumerate}

\end{multicols}
\end{document}
